\section{Introduction}
\label{sec:intro}

A Lichess puzzle earns a Glicko-2 rating in exactly the same way a player does: each user attempt is scored as a win or loss, and the puzzle's rating drifts toward whatever population of solvers consistently succeeds. This makes puzzle ratings a rare well-calibrated continuous difficulty signal on a 6M-position dataset. It also means the ratings take hundreds of attempts to stabilize --- a cheap ex-ante difficulty estimator that reads only the position and solution would let puzzle curators, tutors, and adaptive-training UIs rank a fresh puzzle before any user data exists.

The peer-reviewed literature on puzzle-rating regression has bifurcated. On the transformer side, \glickformer~\citep{milosz2024glickformer} reaches MAE 217.7 on 4.2M puzzles with a factorized spatio-temporal encoder over the sequence of positions, but requires multiple RTX 6000 GPUs and multi-day training. On the shallow side, product blogs (\eg Aimchess's ``median rating by move count'') give simple lookup tables with error near MAE 380. The IEEE BigData 2024 Cup~\citep{ieee2024cup} sat between the two, with the third-place entry combining \maia features~\citep{mcilroy2020aligning} and \lightgbm; details of that stack are not fully public.

{\bf What is missing is a compact, reproducible, CPU-only gradient-boosting reference} that lets practitioners score any Lichess puzzle in milliseconds using only \pychess and scikit-learn, and that quantifies how much of the SOTA gap is inherent noise versus missing representation. Existing feature-engineered baselines under-explore \emph{tactical density}: how many candidate moves a solver faces at each ply, how many are captures, how many are checks. A puzzle whose 8-ply solution has 30 candidate moves per side is dramatically harder than one whose 8-ply solution has 3 candidates per side, yet neither the number of moves nor a mate-flag captures this signal.

We propose \ours, a 155-dimensional feature-engineered histogram gradient-boosted regressor. \ours augments a 141-dimensional parent feature vector (static piece counts, solution-sequence descriptors, and multi-hot theme flags) with a 14-feature \emph{tactical-density block}: per-ply legal-move breadth (\texttt{sol\_legal\_\{max, mean, last, sum\}}), per-ply legal-capture density, per-ply legal-check density, and four setup-move descriptors (\texttt{setup\_is\_capture}, \texttt{setup\_gives\_check}, \texttt{setup\_piece\_type}, \texttt{setup\_to\_center}). We combine \ours with a convex blend of isotonic and CDF quantile-matching calibrators fit on a held-out validation slice; the blend weight is selected by argmin over a mid-band calibration proxy under an explicit val-MAE budget, ensuring no regression is possible.

On a 50{,}000-puzzle hash-bucket test slice, \ours reaches test MAE 240.8 \elo, Spearman $\rho = 0.768$, and middle-decile calibration 85.6 \elo. This beats the global-mean baseline by 148.8 \elo (38\%), a per-theme-mean baseline by 93.1 \elo (28\%), and the 141-feature parent pipeline by 12.8 \elo (5.0\%). The middle-decile calibration figure clears the \textless\,100 \elo target for the first time on this pipeline. Cross-seed reproducibility is essentially exact ($\Delta\text{MAE} = 0.12$), and both seeds independently pick calibrator blend weight $\alpha = 0.30$.

Our contributions are:
\begin{itemize}[leftmargin=*,itemsep=0pt,topsep=0pt]
    \item We propose a 14-feature tactical-density block --- per-ply legal-move, legal-capture, and legal-check counts along the solution, plus setup-move descriptors --- that directly encodes candidate-move ambiguity and forcing structure. On introduction, three of the new features enter the top-15 correlation-proxy importance list.
    \item We propose a budgeted, convex isotonic--CDF calibrator blend whose blend weight $\alpha$ is chosen by argmin over a mid-band calibration proxy under an explicit val-MAE cushion. The procedure is stable across seeds and preserves Spearman rank order exactly.
    \item We conduct a controlled ablation against a 141-feature parent pipeline and show that the entire 12.8 \elo test-MAE improvement, and the 15.3 \elo middle-decile calibration improvement, are attributable to the tactical-density block alone --- calibrator, splits, and hyperparameters are held constant.
    \item We release a CPU-only reference (155-dim feature extraction and blended calibrator) that scores any Lichess puzzle in milliseconds using only \pychess and scikit-learn, and that quantifies the remaining $\approx$23 \elo gap to \glickformer as dominated by inherent Glicko-2 uncertainty at the tails rather than by missing mid-band signal.
\end{itemize}

The rest of the paper is organized as follows. \Secref{sec:related} reviews prior work on puzzle-rating prediction and skill-conditioned move models. \Secref{sec:method} details the dataset splits, feature extraction, booster, and calibrator. \Secref{sec:results} reports headline numbers, per-decile calibration, cross-seed reproducibility, and feature importance. \Secref{sec:discussion} interprets why the tactical-density block closes the middle-decile gap and where the residual gap lies. \Secref{sec:conclusion} summarizes and points to next steps.
